% !TeX root = ../main.tex
\section{方法}\label{sec:method}
令 $y_i$ 为真实值，$\hat y_i$ 为预测值，样本数 $n=6$。
采用式~\eqref{eq:rmse} 的均方根误差，单位与温度相同。
\begin{equation}
 \operatorname{RMSE}=\sqrt{\frac{1}{n}\sum_{i=1}^{n}(\hat y_i-y_i)^2}.
 \label{eq:rmse}
\end{equation}
\begin{align}
 e_i &= \hat y_i-y_i,\label{eq:error}\\
 \operatorname{MSE} &= \frac{1}{n}\sum_{i=1}^{n}e_i^2.\label{eq:mse}
\end{align}
平均绝对误差为 $\operatorname{MAE}=\frac{1}{n}\sum_{i=1}^{n}|e_i|$，
它与 RMSE 一般不同；本组样本中两者恰好相等。
用分段函数标记是否超过容差：
\begin{equation*}
 q_i=\begin{cases}1,&|e_i|\leq 0.5,\\0,&|e_i|>0.5.\end{cases}
\end{equation*}
两个误差样本组成列向量 $\boldsymbol e=\begin{pmatrix}e_1\\e_2\end{pmatrix}$。
